\documentclass[11pt]{article}

\usepackage[margin=1in]{geometry}
\usepackage{amsmath, amssymb, amsthm}
\usepackage{mathtools}
\usepackage{hyperref}

\newtheorem{theorem}{Theorem}
\newtheorem{proposition}[theorem]{Proposition}
\newtheorem{conjecture}[theorem]{Conjecture}
\newtheorem*{theoremA}{Theorem A}
\newtheorem*{theoremAp}{Theorem A$'$}
\newtheorem*{theoremB}{Theorem B}
\newtheorem*{theoremC}{Theorem C}
\newtheorem*{theoremD}{Theorem D}
\newtheorem*{conjectureC}{Conjecture C}

\DeclareMathOperator{\SL}{SL}
\DeclareMathOperator{\Re}{Re}
\DeclareMathOperator{\Im}{Im}

\title{P12. The $\chi_4$ bilinear identity for $\SL_2(\mathbb{N}_0)$ \\
and the Landau twisted divisor sum}
\author{}
\date{}

\begin{document}
\maketitle

\begin{quote}
\textbf{Pointwise theorem (proved).} For every
$A = (a,b,c,d) \in \SL_2(\mathbb{N}_0)$,
\[
\chi_4(a-b) \cdot \chi_4(c+d) \;=\; \chi_4(n+1), \qquad n = ac+bd = \chi(A),
\]
where $\chi_4$ is the non-trivial Dirichlet character mod $4$. Both sides are
simultaneously zero (when $n$ is odd) or both in $\{\pm 1\}$.

\textbf{Consequence.} The bilinear character sum
$\sum_{A:\,\chi(A) \le N} \chi_4(a-b)\,\chi_4(c+d)$ collapses to the twisted
divisor sum
\[
T(N) \;=\; \sum_{n \le N} \tau(n^2+1)\,\chi_4(n+1),
\]
which is empirically $O(\sqrt N)$ over $N \in [10^3, 10^7]$. The cancellation
comes from $\chi_4(n+1)$ averaging over even $n$ where $\tau(n^2+1)$ has positive
density.
\end{quote}

\section*{A clarification (correcting the earlier draft)}

The character used in the identity below is \emph{not} the multiplicative Hecke
character $\sigma$ on $\mathbb{Z}[i]$ that an earlier draft of this note
advertised. Rather, it is the simpler function
$\rho(\alpha) := \chi_4\bigl(\Re(\alpha) + \Im(\alpha)\bigr)$, which is
\emph{not} multiplicative on $\mathbb{Z}[i]$. The bilinear identity below is
\emph{not} a consequence of multiplicativity --- it is a consequence of the
\textbf{determinant constraint} $ad - bc = 1$. The proof exploits a parity
calculation on $bd$ using the determinant.

(For comparison: a truly multiplicative quadratic Hecke character $\sigma$
would give $\sigma(\xi)\sigma(\eta) = \sigma(\xi\eta) = \sigma(n+i)$ trivially,
but $\sigma(n+i)$ is constant in sign on the slice in a way that gives \emph{no}
cancellation in the cumulative sum. So the multiplicative path actually fails.
The non-multiplicative $\rho$ succeeds because its ``non-multiplicativity
defect'' is killed by the $\SL_2$ determinant constraint.)

\section{Definitions}

$\chi_4 : \mathbb{Z} \to \{0, \pm 1\}$ is the unique non-trivial Dirichlet
character mod $4$:
\[
\chi_4(m) \;=\;
\begin{cases}
+1 & m \equiv 1 \pmod 4, \\
-1 & m \equiv 3 \pmod 4, \\
0 & m \text{ even}.
\end{cases}
\]
For $A = \begin{pmatrix} a & b \\ c & d \end{pmatrix} \in \SL_2(\mathbb{N}_0)$,
define
\[
\rho(A) \;:=\; \chi_4(a-b) \cdot \chi_4(c+d).
\]
By Diophantus, $\xi\eta = n + i$ where $\xi = a - bi$, $\eta = c + di$,
$n = \chi(A) = ac + bd$.

\section{Theorem A (Pointwise identity)}

\begin{theoremA}
For every $A \in \SL_2(\mathbb{N}_0)$ with $\det A = 1$ and entries $\ge 0$,
\[
\rho(A) \;=\; \chi_4(n+1).
\]
\end{theoremA}

\begin{proof}
\textbf{Step 1: Parity case analysis.} From $ad - bc = 1$, exactly one of
$\{ad, bc\}$ is odd modulo $2$.
\begin{itemize}
\item \emph{Case A:} $ad$ odd, $bc$ even. Then $a, d$ are both odd, and at least
one of $b, c$ is even.
\item \emph{Case B:} $bc$ odd, $ad$ even. Then $b, c$ are both odd, and at least
one of $a, d$ is even.
\end{itemize}

\noindent\textbf{Step 2: When is $\rho(A) = 0$?}

$\rho(A) \ne 0$ iff $\chi_4(a-b) \ne 0$ and $\chi_4(c+d) \ne 0$ iff $a-b$ odd
and $c+d$ odd iff $\{a,b\}$ have different parity and $\{c,d\}$ have different
parity.
\begin{itemize}
\item In Case A ($a, d$ odd, at least one of $b, c$ even): $a - b$ odd iff $b$
even; $c + d$ odd iff $c$ even. So $\rho(A) \ne 0$ in Case A iff
\textbf{$b, c$ both even}.
\item In Case B ($b, c$ odd, at least one of $a, d$ even): $a - b$ odd iff $a$
even; $c + d$ odd iff $d$ even. So $\rho(A) \ne 0$ in Case B iff
\textbf{$a, d$ both even}.
\end{itemize}

\noindent\textbf{Step 3: Show $\chi_4(n+1) = 0$ when $\rho(A) = 0$.}

$\chi_4(n+1) = 0 \iff n$ odd.

\emph{Sub-case A1} ($a, d$ odd, exactly one of $\{b,c\}$ odd): WLOG $b$ odd,
$c$ even. Then $n = ac + bd = (\text{odd})(\text{even}) +
(\text{odd})(\text{odd}) = \text{even} + \text{odd} = \text{odd}$. \checkmark

\emph{Sub-case A2} ($a, d$ odd, $b$ even, $c$ odd):
$n = (\text{odd})(\text{odd}) + (\text{even})(\text{odd}) = \text{odd} +
\text{even} = \text{odd}$. \checkmark

\emph{Sub-case B1} ($b, c$ odd, $a$ odd, $d$ even):
$n = (\text{odd})(\text{odd}) + (\text{odd})(\text{even}) = \text{odd}$.
\checkmark

\emph{Sub-case B2} ($b, c$ odd, $a$ even, $d$ odd):
$n = (\text{even})(\text{odd}) + (\text{odd})(\text{odd}) = \text{odd}$.
\checkmark

In all cases where $\rho(A) = 0$, $n$ is odd, hence $\chi_4(n+1) = 0$.

\medskip
\noindent\textbf{Step 4: The identity when $\rho(A) \ne 0$.}

Either Case A with $b, c$ both even, or Case B with $a, d$ both even. In
both cases, \textbf{$bd$ is even} (Case A: $b$ even; Case B: $d$ even).

Since $a-b$ and $c+d$ are both odd,
\[
\chi_4(a-b)\,\chi_4(c+d) \;=\; \chi_4\bigl((a-b)(c+d)\bigr)
\]
(multiplicativity of $\chi_4$ holds when both arguments are odd, which they
are). Now compute
\begin{align*}
(a-b)(c+d) - (n+1) &= (ac + ad - bc - bd) - (ac + bd + 1) \\
&= ad - bc - 2bd - 1 \\
&= (ad - bc) - 2bd - 1.
\end{align*}
Using $ad - bc = 1$,
\[
(a-b)(c+d) - (n+1) \;=\; 1 - 2bd - 1 \;=\; -2bd.
\]
Since $bd$ is even, $2bd \equiv 0 \pmod 4$, so
\[
(a-b)(c+d) \equiv n+1 \pmod 4.
\]
Both sides are odd. As $\chi_4$ depends only on the residue mod $4$,
$\chi_4\bigl((a-b)(c+d)\bigr) = \chi_4(n+1)$.
\end{proof}

\section{Theorem B (Bilinear-to-divisor reduction)}

\begin{theoremB}
Let $T(N) := \sum_{A \in \SL_2(\mathbb{N}_0),\, A \ne I:\, \chi(A) \le N} \rho(A)$.
Then
\[
T(N) \;=\; \sum_{n=1}^{N} \tau(n^2+1)\,\chi_4(n+1).
\]
\end{theoremB}

(Equivalently, $\sum_{A:\,\chi(A) \le N} \rho(A) = 1 + T(N)$, the $1$
contributed by the identity matrix $I$ at $n=0$, where
$\rho(I) = \chi_4(1)\chi_4(1) = 1$.)

\begin{proof}
By Theorem A, $\rho(A) = \chi_4(\chi(A)+1)$ depends only on $n = \chi(A)$.
Group the sum by $n$:
\[
T(N) + \rho(I) \;=\; \sum_{n=0}^{N} \chi_4(n+1) \cdot
\#\{A \in \SL_2(\mathbb{N}_0) : \chi(A) = n\}.
\]
By the Shakov bijection
$\widehat F_{\phi_0} : \SL_2(\mathbb{N}_0) \xrightarrow{\sim}
\mathcal{D}_{\phi_0}$, the set $\{A : \chi(A) = n\}$ is in bijection with
divisor pairs $(m, n)$, $m \mid n^2+1$. There are $\tau(n^2+1)$ such pairs.
For $n = 0$: $\tau(1) = 1$ (the identity $I$ alone) and $\chi_4(1) = 1$.
Subtract this term to obtain the stated equation.
\end{proof}

\section{Cancellation in $T(N)$ --- what is now proved}

\subsection*{Theorem C (unconditional)}

\begin{theoremC}
$T(N) = O(N)$ unconditionally. Explicitly,
\[
|T(N)| \;\le\; 2 \sum_{\substack{d \in L_{\mathrm{odd}} \\ 1 \le d \le N}}
2^{\omega(d)},
\]
where $L_{\mathrm{odd}}$ is the set of odd integers all of whose prime factors
are $\equiv 1 \pmod 4$, and $\omega(d)$ is the number of distinct prime
factors of $d$.
\end{theoremC}

This RHS is asymptotic to $\sim c \cdot N$ for an explicit positive constant $c$
(with $2c \approx 0.637$ confirmed numerically). The asymptotic constant
depends on the effective form of Selberg--Delange; the uniform bound
$|T(N)| = O(N)$ is what we prove.

This is a $\log N$ improvement over the trivial bound
$|T(N)| \le \sum_{n \le N} \tau(n^2+1) \asymp \tfrac{3}{\pi} N \log N$.

\begin{proof}
\textbf{Step 1 (Hyperbola).} For $n \ge 1$, $n^2+1$ is never a perfect square
(since $n^2 < n^2+1 < (n+1)^2$ for $n \ge 1$). Hence
$\tau(n^2+1) = 2 \cdot \#\{d : d \mid n^2+1,\, 1 \le d \le n\}$. So
\[
T(N) \;=\; 2 \sum_{n=1}^N \chi_4(n+1) \sum_{\substack{d \mid n^2+1 \\ 1 \le d \le n}} 1
\;=\; 2 \sum_{d=1}^N \sum_{\substack{n:\, d \mid n^2+1 \\ d \le n \le N}} \chi_4(n+1).
\]

\noindent\textbf{Step 2 (Restriction on $d$).} Define
$\rho(d) := \#\{x \pmod d : x^2 \equiv -1 \pmod d\}$. Classically,
$\rho(d) > 0$ iff $v_2(d) \le 1$ and every odd prime factor of $d$ is
$\equiv 1 \pmod 4$. When $\rho(d) > 0$,
$\rho(d) = 2^{\omega(d_{\mathrm{odd}})}$ where $d_{\mathrm{odd}}$ is the odd
part of $d$ (CRT: at $p = 2$, $\rho(2) = 1$, $\rho(2^k) = 0$ for $k \ge 2$;
at odd $p \equiv 1 \pmod 4$, $\rho(p^k) = 2$ for $k \ge 1$).

If $\rho(d) = 0$, the inner set is empty, contributing $0$. So we restrict to
$\rho(d) > 0$.

\medskip
\noindent\textbf{Step 3 (Even $d$ contribute zero).} Suppose $d = 2 d'$ with
$d'$ odd. Then $\rho(d) > 0$ requires $d' \in L_{\mathrm{odd}}$, and
$\rho(d) = \rho(2)\rho(d') = 1 \cdot 2^{\omega(d')}$. Crucially, $\rho(2) = 1$
corresponds to $x \equiv 1 \pmod 2$, i.e.\ $x$ odd. So every $n_0 \pmod d$
with $n_0^2 \equiv -1 \pmod d$ has $n_0$ odd. Hence $n \equiv n_0 \pmod d$
implies $n$ odd, hence $n+1$ even, hence $\chi_4(n+1) = 0$. So all such $d$
contribute $0$ to $T(N)$.

Therefore the only contributing $d$ are those in $L_{\mathrm{odd}}$ ($d$ odd,
all prime factors $\equiv 1 \pmod 4$).

\medskip
\noindent\textbf{Step 4 (Inner AP-sum bound).} Fix $d \in L_{\mathrm{odd}}$
with $d \le N$, and $n_0 \in \{0, 1, \dots, d-1\}$ with
$n_0^2 \equiv -1 \pmod d$. Consider
\[
S(d, n_0) \;:=\; \sum_{\substack{n:\, n \equiv n_0 \pmod d \\ d \le n \le N}} \chi_4(n+1).
\]
Writing $n = n_0 + j d$, the constraint $d \le n \le N$ gives $j$ in an
interval of length at most $\lfloor (N - n_0)/d \rfloor + 1$.

As $j$ varies through any $4$ consecutive integers, $n_0 + 1 + jd \pmod 4$
runs through all $4$ residues mod $4$ (because $\gcd(d, 4) = 1$). The sum of
$\chi_4$ over a complete set of residues mod $4$ is
$\chi_4(0) + \chi_4(1) + \chi_4(2) + \chi_4(3) = 0 + 1 + 0 + (-1) = 0$.

Hence $S(d, n_0)$ equals (sum over remainder of $\le 3$ extra $j$-values), which
has absolute value at most $1$ (partial sums of $\chi_4$ on an AP with step
coprime to $4$ are bounded by $1$).

Therefore $|S(d, n_0)| \le 1$.

\medskip
\noindent\textbf{Step 5 (Total bound).}
\begin{align*}
|T(N)| &\le 2 \sum_{\substack{d \in L_{\mathrm{odd}} \\ d \le N}}
\sum_{\substack{n_0 \pmod d \\ n_0^2 \equiv -1 \pmod d}} |S(d, n_0)| \\
&\le 2 \sum_{\substack{d \in L_{\mathrm{odd}} \\ d \le N}} \rho(d) \cdot 1 \\
&= 2 \sum_{\substack{d \in L_{\mathrm{odd}} \\ d \le N}} 2^{\omega(d)}.
\end{align*}

\medskip
\noindent\textbf{Step 6 (Selberg--Delange).} The multiplicative function
$f(d) := 2^{\omega(d)} \mathbf{1}[d \in L_{\mathrm{odd}}]$ has Dirichlet series
\[
F(s) \;=\; \sum_{d \ge 1} \frac{f(d)}{d^s}
\;=\; \prod_{p \equiv 1 \!\!\!\pmod 4} \frac{1 + p^{-s}}{1 - p^{-s}}.
\]
Each Euler factor has a pole of order $1$ at $s = 1$ on primes
$p \equiv 1 \pmod 4$, which have density $1/2$ among primes. By
Selberg--Delange (Tenenbaum, \emph{Introduction to Analytic and Probabilistic
Number Theory}, Ch.~II.5, Thm.~5.2), for a multiplicative $f$ with
$\sum_{p \le x} f(p)/p = \kappa \log\log x + O(1)$ and bounded $f(p^k)$,
\[
\sum_{d \le X} f(d) \;\sim\; c\,X (\log X)^{\kappa - 1}
\quad (X \to \infty).
\]
Here
$\sum_{p \le X} f(p)/p = 2 \sum_{\substack{p \le X \\ p \equiv 1 \!\pmod 4}}
1/p = 2 \cdot (1/2) \log \log X + O(1) = \log \log X + O(1)$,
giving $\kappa = 1$. Therefore
\[
\sum_{\substack{d \le N \\ d \in L_{\mathrm{odd}}}} 2^{\omega(d)} \;\sim\; c \cdot N
\]
for an explicit $c > 0$.

Combining Steps 5 and 6,
$|T(N)| \le 2 \sum_{d \le N,\, d \in L_{\mathrm{odd}}} 2^{\omega(d)} \sim 2cN$,
hence $|T(N)| = O(N)$ unconditionally.
\end{proof}

\subsection*{Numerical confirmation}

The constant in the bound is $\approx 0.637$:

\begin{center}
\begin{tabular}{r|r|r|r}
$N$ & bound $2 \sum 2^{\omega(d)}$ & bound$/N$ & actual $|T(N)|$ \\ \hline
$100$ & $66$ & $0.660$ & $2$ \\
$1000$ & $634$ & $0.634$ & $16$ \\
$10000$ & $6374$ & $0.637$ & $6$ \\
$50000$ & $31842$ & $0.637$ & $112$
\end{tabular}
\end{center}

The bound is loose by a factor of $\sim \sqrt N / \log N$ compared to the
empirical $|T(N)| \sim \sqrt N$, but is rigorous and unconditional.

\subsection*{Conjecture C (open)}

\begin{conjectureC}
$|T(N)| = O\bigl(\sqrt N\,(\log N)^{O(1)}\bigr)$.
\end{conjectureC}

Empirically $|T(N)|/\sqrt N \in [0.06, 1.45]$ over
$N \in \{10^3, \dots, 10^7\}$. Closing the gap from $O(N)$ to $O(\sqrt N)$
would require non-trivial analytic NT (subconvexity for an appropriate Hecke
$L$-function over $\mathbb{Q}(i)$, or shifted-convolution machinery).

\section{What's actually proved here}

\textbf{Proved unconditionally:}
\begin{enumerate}
\item Theorem A --- the pointwise identity. Elementary, parity calculation on
$bd$ via the determinant.
\item Theorem B --- reduction to a twisted divisor sum, via the Shakov bijection.
\item Theorem C --- $T(N) = O(N)$ unconditionally by elementary AP-sum
cancellation. \emph{Strictly weaker than the empirical $O(\sqrt N)$.}
\end{enumerate}

\textbf{Conjecture C:} $T(N) = O(\sqrt N (\log N)^{O(1)})$. Empirically supported
across $N \in [10^3, 10^7]$. Would follow from subconvexity for an appropriate
Hecke $L$-function on $\mathbb{Q}(i)$.

\section{A second linear-form identity, and the limitation}

By exactly the same parity argument as Theorem A:

\begin{theoremAp}
$\chi_4(a + b) \cdot \chi_4(c - d) = \chi_4(n - 1)$ for every
$A \in \SL_2(\mathbb{N}_0)$.
\end{theoremAp}

\begin{proof}
$(a+b)(c-d) - (n-1) = (ac - ad + bc - bd) - (ac + bd - 1) = -ad + bc - 2bd + 1
= -(ad-bc) - 2bd + 1 = -1 - 2bd + 1 = -2bd$. The same determinant-driven
parity analysis kills $-2bd$ mod $4$ when
$\rho_2(A) := \chi_4(a+b)\chi_4(c-d) \ne 0$.
\end{proof}

\noindent For $n$ even, $\chi_4(n-1) = -\chi_4(n+1)$ (verified directly). So
$\rho_2(A) = -\rho(A)$ --- the second identity gives the \emph{same}
information up to sign.

\subsection*{Characterization of all such linear-form identities}

(This proof was rewritten in round~2 after a referee pointed out that the earlier
``exact-coefficient-matching'' argument was wrong: the canonical example
$(a-b)(c+d) = ac - bd + 1$ has $bd$-coefficient $-1$, not $+1$, so exact
matching fails. The correct mechanism uses \textbf{mod-4 matching with the
parity constraint on $bd, ac, bc$} that follows from $A \in \SL_2(\mathbb{N}_0)$
with $L_1, L_2$ both odd.)

\begin{theoremD}
Suppose $L_1(a, b) = \alpha a + \beta b$ and $L_2(c, d) = \gamma c + \delta d$
are integer linear forms such that
\[
\chi_4(L_1(a,b)) \cdot \chi_4(L_2(c,d)) = \chi_4(c_1 n + c_0) \quad
\text{for all } A \in \SL_2(\mathbb{N}_0)
\]
for some constants $c_0, c_1$ (in $\mathbb{Z}/4\mathbb{Z}$). Then, viewing
$(\alpha, \beta, \gamma, \delta) \in (\mathbb{Z}/4\mathbb{Z})^4$, the
solutions are exactly the pairs equivalent (mod $4$, up to overall sign) to
one of
\[
(L_1, L_2) \in \{(a - b,\, c + d),\, (a + b,\, c - d)\}.
\]
The two yield $\chi_4(n+1)$ and $\chi_4(n-1)$ respectively. Since
$\chi_4(n-1) = -\chi_4(n+1)$ for $n$ even (the only case both characters are
non-zero), the two identities are equivalent up to sign.
\end{theoremD}

\begin{proof}
By Theorem A's case analysis, $\chi_4(L_1)\chi_4(L_2) \ne 0$ iff both $L_1, L_2$
are odd, which restricts $A$ to:
\begin{itemize}
\item \textbf{Case A}: $a, d$ odd, $b, c$ both even; then $ac \in 2\mathbb{Z}$,
$bd \in 2\mathbb{Z}$, $bc \in 4\mathbb{Z}$.
\item \textbf{Case B}: $b, c$ odd, $a, d$ both even; then $ac \in 2\mathbb{Z}$,
$bd \in 2\mathbb{Z}$, $bc$ odd.
\end{itemize}
(In Case A, $a$ odd and $c$ even gives $ac$ even with $ac \pmod 4 \in \{0, 2\}$;
similarly $bd$. In Case B, both even-times-odd, same. And $bc$ is in
$4\mathbb{Z}$ in Case A vs.\ odd in Case B.)

Expand and substitute $ad = 1 + bc$:
\[
L_1 L_2 = \alpha\gamma \cdot ac + \beta\delta \cdot bd
+ (\alpha\delta + \beta\gamma) \cdot bc + \alpha\delta.
\]
We require $L_1 L_2 \equiv c_1(ac + bd) + c_0 \pmod 4$ for all $(a,b,c,d)$ in
the relevant cases. Equivalently:
\[
\underbrace{(\alpha\gamma - c_1)\,ac}_{(*)} +
\underbrace{(\beta\delta - c_1)\,bd}_{(**)} +
\underbrace{(\alpha\delta + \beta\gamma)\,bc}_{(\dagger)} +
\underbrace{(\alpha\delta - c_0)}_{(\ddagger)} \equiv 0 \pmod 4.
\]

Each term is extracted in two stages: first using Case A, where $bc \in 4\mathbb{Z}$
kills $(\dagger)$, then using Case B, where $bc$ is odd. Within each case, $ac$,
$bd$, $bc$ vary freely over their allowed residues mod $4$ as $A$ ranges over
$\SL_2(\mathbb{N}_0)$ matrices in that case (computationally verified: e.g.,
$S^2 T = (3,2,1,1)$ has $ac=3$, $bd=2$; $T^2 S = (1,1,2,3)$ has $ac=2$, $bd=3$;
etc.).

\smallskip
\textbf{Case A first} ($bc \in 4\mathbb{Z}$, so $(\dagger) \equiv 0 \pmod 4$
automatically):
\begin{itemize}
\item $(*)$: $ac$ takes both values mod $4$ in $\{0, 2\}$ (e.g., $a=1, c=2$
gives $ac=2$; $a=1, c=4$ gives $ac=4 \equiv 0$). Hence
$(\alpha\gamma - c_1) \cdot 2 \equiv 0 \pmod 4$ $\Rightarrow$
$\alpha\gamma \equiv c_1 \pmod 2$.
\item $(**)$: same argument, $\beta\delta \equiv c_1 \pmod 2$.
\item $(\ddagger)$: $\alpha\delta \equiv c_0 \pmod 4$ (extracted as the constant
after $(*)$, $(**)$, $(\dagger)$ are accounted for).
\end{itemize}

\textbf{Then Case B} ($bc$ odd): now $(*)$ and $(**)$ already hold, so the
residual constraint is $(\dagger) + (\ddagger) \equiv 0 \pmod 4$, i.e.,
$(\alpha\delta + \beta\gamma) \cdot bc + (\alpha\delta - c_0) \equiv 0 \pmod 4$.
Subtracting the already-established $(\ddagger)$ leaves
$(\alpha\delta + \beta\gamma) \cdot bc \equiv 0 \pmod 4$ for $bc$ ranging over
odd integers. Since (odd) $\times k \equiv 0 \pmod 4$ forces
$k \equiv 0 \pmod 4$:
\begin{itemize}
\item $(\dagger)$: $\alpha\delta + \beta\gamma \equiv 0 \pmod 4$.
\end{itemize}

For a non-degenerate (i.e., $\chi_4 \ne 0$) identity, we need $\alpha\gamma$
odd, equivalently both $\alpha, \gamma$ odd. Similarly $\beta\delta$ odd, so
$\beta, \delta$ odd. So all four coefficients are in $\{1, -1\} \pmod 4$.

Then $\alpha\delta, \beta\gamma \in \{1, -1\} \pmod 4$, and $(\dagger)$
requires $\alpha\delta + \beta\gamma \equiv 0 \pmod 4$. Since each is $\pm 1$,
the only way to sum to $0 \pmod 4$ is $\{+1, -1\}$, i.e.,
$\alpha\delta = -\beta\gamma$. Two sub-cases:
\begin{itemize}
\item[(a)] $\alpha\delta = 1$, $\beta\gamma = -1$: forces $\alpha \equiv \delta$,
$\beta \equiv -\gamma$ (both mod $4$). Up to overall sign in $L_1$ and $L_2$,
this reduces to $(L_1, L_2) = (a-b, c+d)$.
\item[(b)] $\alpha\delta = -1$, $\beta\gamma = 1$: $\alpha \equiv -\delta$,
$\beta \equiv \gamma$. Reduces to $(L_1, L_2) = (a+b, c-d)$ up to signs.
\end{itemize}
In both cases, the resulting identity is $\chi_4(n+1)$ (case (a), with
$c_0 = \alpha\delta = 1$) or $\chi_4(n-1)$ (case (b), with
$c_0 = \alpha\delta = -1$). A sign change in $c_1 = \alpha\gamma$
corresponds to swapping $L_1 \to -L_1$ or $L_2 \to -L_2$, which gives
equivalent characters.
\end{proof}

\paragraph{Consequence: the linear-form Plancherel program is sharply limited
at conductor 4.}
Higher conductor (e.g., $\chi_8$) would require the defect $-2bd$ to be
$\equiv 0 \pmod 8$, i.e., $bd \equiv 0 \pmod 4$ --- but in $\SL_2(\mathbb{N}_0)$
we only have $bd$ even (and $bd \equiv 2 \pmod 4$ does occur, as in
$A = (1, 2, 0, 1)$). So the Plancherel program \textbf{does not extend} to
higher-conductor linear-form identities by this technique. New ideas are
needed.

\section{What this gives toward Landau IV}

\paragraph{The honest scope.}
We have ONE identity at conductor $4$, giving square-root cancellation in the
corresponding twisted divisor sum \emph{empirically}. The
Friedlander--Iwaniec asymptotic sieve requires bilinear estimates for
\emph{all} Type-II inputs.

\paragraph{What remains for Landau IV via FI.}
Decompose general $\alpha(\xi)\beta(\eta)$ Plancherel-style into:
\begin{itemize}
\item Linear-form characters at conductor $4$: covered by Theorem A (one
identity, up to sign).
\item Linear-form characters at higher conductors: NOT covered by Theorem A's
technique (Theorem D's limitation).
\item Multiplicative Hecke characters on $\mathbb{Z}[i]$: angular characters
degenerate (P3); a true multiplicative quadratic Hecke character $\sigma$ gives
constant sign on the slice and no cancellation.
\item Mixed / non-decomposable inputs.
\end{itemize}

\paragraph{The honest open problem.}
Find a NEW class of identities (probably involving quadratic forms in $a, b,
c, d$, or higher-order constructions, or arithmetic data of $A$ beyond the
linear coordinates) that captures the cancellation needed for the
higher-conductor part of the Plancherel decomposition. Theorem A is one piece
--- likely a small piece --- of this larger structure.

\section{What is NEW and what is CLASSICAL}

\textbf{New (this note):}
\begin{itemize}
\item The \textbf{pointwise identity} $\chi_4(a-b)\chi_4(c+d) = \chi_4(n+1)$ on
$\SL_2(\mathbb{N}_0)$ --- a consequence of the determinant constraint, not
previously stated in the project.
\item The clean \textbf{collapse} of the bilinear $\rho$-sum to the twisted
divisor sum $T(N)$.
\item The identification of $\rho$ as the \emph{first non-degenerate} spin
candidate in the Shakov framework (in contrast to P1's word-parity, P3's
angular).
\end{itemize}

\textbf{Classical / external:}
\begin{itemize}
\item The bound $T(N) = O(N^{1/2+\varepsilon})$ for the twisted divisor sum is
in the analytic NT literature (Hooley, Iwaniec). Our contribution is
\emph{identifying} that the bilinear sum reduces to this object.
\end{itemize}

\textbf{Plancherel completeness (open):} whether \emph{every} Type-II input has
this kind of clean collapse via a determinant-constraint-driven identity.

\section{Open problems}

\paragraph{Q12.1.} Write down the precise Hecke $L$-function whose values give
$T(N)$. (Likely $L(s, \chi_4 \otimes \theta_4)$ where $\theta_4$ is the binary
theta of $x^2+y^2$.)

\paragraph{Q12.2.} Find a $(\chi, \chi')$ --- both Dirichlet characters of small
conductor on linear functionals of $(a, b)$, $(c, d)$ --- giving an analogous
pointwise identity. Combinations to test:
\begin{itemize}
\item $\chi_4(a+b)\chi_4(c-d)$
\item $\chi_q(a)\chi_q(c)$ for various $q$
\item $\chi_q(b)\chi_q(d)$ for various $q$
\item Mixed pairings at different conductors.
\end{itemize}

\paragraph{Q12.3.} Identify the family of $(\chi, \chi')$ pairs for which a
pointwise identity $\chi(L_1(a,b))\chi'(L_2(c,d)) = \mu(\chi, \chi'; n)$ holds
(with $\mu$ a divisor-like function of $n$), via a determinant-driven parity
calculation analogous to Theorem A's Step 4.

\paragraph{Q12.4.} Decompose a generic FI Type-II input $\alpha \otimes \beta$
along this family, and aggregate the cancellations.

\paragraph{Q12.5.} Generalize to $\phi_1, \psi_2, \phi_3$. Each has its own
determinant constraint and its own analogue of $\rho$.

\section{Summary}

The \textbf{proved core result} (Theorems A and B): a pointwise determinant-driven
identity collapses a bilinear character sum on $\SL_2(\mathbb{N}_0)$ to a
twisted divisor sum $T(N) = \sum_{n \le N} \tau(n^2+1)\,\chi_4(n+1)$. This
sum has $O(N^{1/2+\varepsilon})$ cancellation by classical methods
(Hooley-style).

This is the \textbf{first power-saving spin} in the Shakov framework, succeeding
where word-parity (P1) gives linear positive bias and angular characters (P3)
degenerate on the slice. It opens a Plancherel-style program to recover the
full FI bilinear bound from a family of such identities --- the precise next
research target.

\end{document}
